\documentclass[11pt]{article}
\usepackage[margin=1in]{geometry}
\usepackage{amsmath,amssymb,amsthm}

\newtheorem{theorem}{Theorem}
\newtheorem{lemma}[theorem]{Lemma}
\DeclareMathOperator{\ch}{ch}

\begin{document}

\title{Irreducible Vertices in Positive-definite Tree Lattices: Solution}
\date{}
\maketitle

We use $x^2$ for $x\mathbin{\cdot}x$.  If $C$ is a weighted tree, then
$Q_C$ denotes its Gram matrix in the vertex basis.  Coefficients of a lattice
element $x=\sum_y x_y y$ are denoted by $x_y$.

\begin{lemma}[Norm-one elements]
If $L$ is a positive-definite integral lattice and $x\in L$ satisfies
$x^2=1$, then $x$ is irreducible.
\end{lemma}

\begin{proof}
If $x=a+b$ with $a,b\neq 0$ and $a\mathbin{\cdot}b\geq0$, integrality and
positive definiteness give $a^2,b^2\in\mathbb Z_{>0}$.  Hence
\[
 1=x^2=a^2+b^2+2a\mathbin{\cdot}b\geq2,
\]
a contradiction.
\end{proof}

Root a weighted tree $C$ at $\rho$ and orient its edges away from $\rho$.
Let $\ch_C(y)$ be the number of children of $y$.  Call $C$ \emph{admissible}
when its form is positive definite, all its weights are at least $2$, and
\[
 w(y)\geq \ch_C(y)+1\qquad(y\in C).
\]
For a positive-definite rooted tree define its \emph{capacity} by
\[
 \gamma(C)=(Q_C^{-1})_{\rho\rho}.
\]
The inverse exists because $Q_C$ is positive definite.  It has rational
entries, so $\gamma(C)\in\mathbb Q$; inequalities below may equivalently be
read after extending scalars to $\mathbb R$.

\begin{lemma}[Capacity recursion]
Let $C$ be admissible.  If the root has children
$\rho_1,\ldots,\rho_t$ and $C_i$ is the rooted subtree at $\rho_i$, then
\[
 \gamma(C)=\frac{1}{w(\rho)-\sum_{i=1}^t\gamma(C_i)}
 \qquad\text{and}\qquad 0<\gamma(C)<1.
\]
\end{lemma}

\begin{proof}
Order the root first and group the other vertices by the subtrees $C_i$.
Then
\[
 Q_C=\begin{pmatrix}w(\rho)&c^T\\c&Q_0\end{pmatrix},
 \qquad Q_0=\bigoplus_i Q_{C_i},
\]
where $c$ has entry $-1$ at each child root and $0$ elsewhere.  Block
Gaussian elimination, or the inverse formula for a Schur complement, gives
\[
 (Q_C^{-1})_{\rho\rho}
 =\frac{1}{w(\rho)-c^TQ_0^{-1}c}
 =\frac{1}{w(\rho)-\sum_i\gamma(C_i)}.
\]
The denominator is positive: it is the Schur complement of the
positive-definite principal block $Q_0$ in the positive-definite matrix
$Q_C$.  This also proves $\gamma(C)>0$.

We prove the upper bound by induction on $|C|$.  For a single vertex,
$\gamma(C)=1/w(\rho)<1$ because $w(\rho)\geq2$.  Otherwise the induction
hypothesis gives $\gamma(C_i)<1$.  Thus, using
$w(\rho)\geq t+1$,
\[
 w(\rho)-\sum_i\gamma(C_i)>w(\rho)-t\geq1,
\]
and the displayed reciprocal formula yields $\gamma(C)<1$.
\end{proof}

\begin{lemma}[Rooted integer estimate]\label{lem:rooted-estimate}
Let $C$ be admissible with root $\rho$ and capacity $\gamma$.  For every
$x\in L(C)$ and every $k\in\mathbb Z$,
\[
 x^2-(2k+1)x_\rho+\gamma k(k+1)\geq0.
\]
\end{lemma}

\begin{proof}
We induct on $|C|$.  If $C$ consists only of $\rho$, write $x=m\rho$ and
$W=w(\rho)$.  Since $\gamma=1/W$, multiplying the asserted inequality by
$W>0$ gives
\[
 W\left(Wm^2-(2k+1)m+\frac{k(k+1)}W\right)
 =(Wm-k)(Wm-k-1)\geq0.
\]
The last inequality holds because the two factors are consecutive integers.

Now let the children of $\rho$ be $\rho_1,\ldots,\rho_t$, let $C_i$ be the
corresponding rooted subtrees, and write
\[
 x=a\rho+\sum_i x_i,\qquad x_i\in L(C_i),\qquad
 s_i=(x_i)_{\rho_i},\qquad\gamma_i=\gamma(C_i).
\]
The tree-lattice form gives
\[
 x^2=w(\rho)a^2-2a\sum_i s_i+\sum_i x_i^2.
\]
Apply the induction hypothesis in $C_i$ first with the integer $a$ and then
with $a-1$.  The two resulting bounds are
\begin{align*}
 x_i^2-2as_i&\geq s_i-\gamma_i a(a+1),\\
 x_i^2-2as_i&\geq-s_i-\gamma_i a(a-1).
\end{align*}
Taking their maximum gives
\[
 x_i^2-2as_i\geq-\gamma_i a^2+|s_i-\gamma_i a|.
\]
Set $D=\sum_i|s_i-\gamma_i a|$.  Since
$\gamma^{-1}=w(\rho)-\sum_i\gamma_i$, and setting $\tau=a/\gamma$, we obtain
\begin{align*}
 x^2-(2k+1)a+\gamma k(k+1)
 &\geq \gamma(\tau-k)(\tau-k-1)+D. \tag{1}
\end{align*}

If $\tau\notin(k,k+1)$, the quadratic term in (1) is nonnegative.  Otherwise
write $\alpha=\tau-k$ and $\beta=k+1-\tau$.  Then
$0<\alpha,\beta<1$, $\alpha+\beta=1$, and the quadratic term is
$-\gamma\alpha\beta$.  It remains to bound $D$.

The integer
\[
 N=w(\rho)a-\sum_i s_i
\]
satisfies
\[
 \tau-N
 =\left(w(\rho)-\sum_i\gamma_i\right)a
   -\left(w(\rho)a-\sum_i s_i\right)
 =\sum_i(s_i-\gamma_i a).
\]
Consequently the triangle inequality gives
\[
 D\geq|\tau-N|\geq\operatorname{dist}(\tau,\mathbb Z)
 =\min\{\alpha,\beta\}.
\]
Because $0<\gamma<1$ and
$\alpha\beta\leq\min\{\alpha,\beta\}$, we have
$D\geq\gamma\alpha\beta$.  Substitution into (1) finishes the induction.
\end{proof}

\begin{lemma}[Root coefficient bound]\label{lem:root-bound}
If $C$ is admissible and $0\neq x\in L(C)$, then
\[
 x^2-x_\rho>0.
\]
\end{lemma}

\begin{proof}
Let $\gamma=\gamma(C)$.  If $x_\rho\leq0$, positive definiteness gives
$x^2-x_\rho>0$.  Suppose $x_\rho>0$.  Let
$\rho^\#\in L(C)\otimes\mathbb Q$ be the dual vector characterized by
\[
 \rho^\#\mathbin{\cdot}y=y_\rho\qquad(y\in L(C)).
\]
Its coefficient vector is $Q_C^{-1}e_\rho$, and therefore
$(\rho^\#)^2=\gamma$.  Cauchy--Schwarz in the positive-definite real
extension gives
\[
 x_\rho^2=(x\mathbin{\cdot}\rho^\#)^2
 \leq x^2(\rho^\#)^2=\gamma x^2.
\]
Since $0<\gamma<1$ and $x_\rho$ is a positive integer,
\[
 x^2\geq\frac{x_\rho^2}{\gamma}>x_\rho^2\geq x_\rho.
\]
\end{proof}

\begin{theorem}
Let $T$ be a finite integer-weighted tree whose tree-lattice form is positive
definite.  If there is a unique vertex $v$ satisfying
$w(v)<d_T(v)$, then $T$ contains a vertex irreducible in $L(T)$.
\end{theorem}

\begin{proof}
Applying positive definiteness to each vertex basis vector shows that every
weight is a positive integer.  If a vertex has weight $1$, it is irreducible
by the norm-one lemma.  We may therefore assume that every weight is at least
$2$; we prove that the unique vertex $v$ with $w(v)<d_T(v)$ is irreducible.

Let $C_1,\ldots,C_m$ be the connected components of $T\setminus\{v\}$.
For each $i$, root $C_i$ at the unique vertex $\rho_i$ adjacent to $v$.
Each $Q_{C_i}$ is a principal submatrix of the positive-definite matrix
$Q_T$, so the form on $C_i$ is positive definite.  Every $x\in C_i$ is not
the exceptional vertex, hence $w(x)\geq d_T(x)$.  With edges oriented away
from $\rho_i$, every such $x$ has one additional incident edge toward its
parent, or toward $v$ when $x=\rho_i$.  Thus
\[
 d_T(x)=\ch_{C_i}(x)+1,
\]
so every $C_i$ is admissible.  Write $\gamma_i=\gamma(C_i)$ and
$W=w(v)$.  The Schur complement at $v$ in $Q_T$ is positive, whence
\[
 W-\sum_i\gamma_i>0. \tag{2}
\]

Assume for a contradiction that $v=a+b$, where $a,b\neq0$ and
$a\mathbin{\cdot}b\geq0$.  Put $z=-b$, so $a=v+z$ and
\[
 z^2+v\mathbin{\cdot}z\leq0. \tag{3}
\]
Moreover, $z\neq0$ and $z\neq-v$.  Let $q$ be the coefficient of $v$ in
$z$.  If $q\leq-1$, replace $z$ by $-v-z$; this exchanges $a$ and $b$,
preserves $z\mathbin{\cdot}(z+v)$, and changes the coefficient to
$-1-q\geq0$.  Thus we may write
\[
 z=pv+\sum_i y_i,qquad p\in\mathbb Z_{\geq0},\qquad y_i\in L(C_i).
\]
Let $s_i=(y_i)_{\rho_i}$.  Since
$v\mathbin{\cdot}y_i=-s_i$, direct expansion gives
\[
 z^2+v\mathbin{\cdot}z
 =Wp(p+1)+\sum_i\bigl(y_i^2-(2p+1)s_i\bigr). \tag{4}
\]

If $p\geq1$, Lemma~\ref{lem:rooted-estimate}, with $k=p$, changes (4) into
\[
 z^2+v\mathbin{\cdot}z
 \geq\left(W-\sum_i\gamma_i\right)p(p+1)>0
\]
by (2), contradicting (3).

If $p=0$, equation (4) becomes
\[
 z^2+v\mathbin{\cdot}z=\sum_i(y_i^2-s_i).
\]
Every summand is zero when $y_i=0$ and strictly positive otherwise by
Lemma~\ref{lem:root-bound}.  At least one $y_i$ is nonzero because $z\neq0$.
The sum is therefore positive, again contradicting (3).  Hence $v$ is
irreducible.
\end{proof}

\end{document}
